\sectioncentered*{Заключение}
\addcontentsline{toc}{section}{Заключение}

Сделана система, которая по PDF определяет, русский это текст или немецкий. Используются короткие слова, частотные слова и нейросеть. Всё запускается локально через Docker Compose.

На четырнадцати проверочных фразах короткие и частотные слова ошибаются на текстах из одних длинных слов. Нейросеть на этой выборке отвечает верно: русский и немецкий различаются алфавитом. На длинном документе быстрее метод частотных слов.

\begin{thebibliography}{9}
	\addcontentsline{toc}{section}{Cписок использованных источников}

	\bibitem{cavnar}
	Cavnar W., Trenkle J. N-Gram-Based Text Categorization // Proceedings of SDAIR-94. \mdash 1994. \mdash Режим доступа: \url{https://www.let.rug.nl/vannoord/TextCat/textcat.pdf}. \mdash Дата доступа: 23.09.2026.

	\bibitem{scikit}
	Pedregosa F. и др. Scikit-learn: Machine Learning in Python // Journal of Machine Learning Research. \mdash 2011. \mdash Vol.~12. \mdash P.~2825--2830.

	\bibitem{fastapi}
	FastAPI documentation [Электронный ресурс]. \mdash Режим доступа: \url{https://fastapi.tiangolo.com/}. \mdash Дата доступа: 23.09.2026.

	\bibitem{pymorphy}
	Korobov M. Morphological Analyzer and Generator for Russian and Ukrainian Languages // Analysis of Images, Social Networks and Texts. \mdash 2015. \mdash P.~320--332.

	\bibitem{pypdf}
	pypdf [Электронный ресурс]. \mdash Режим доступа: \url{https://pypdf.readthedocs.io/}. \mdash Дата доступа: 23.09.2026.

	\bibitem{postgres}
	PostgreSQL 16 Documentation [Электронный ресурс]. \mdash Режим доступа: \url{https://www.postgresql.org/docs/16/}. \mdash Дата доступа: 23.09.2026.

	\bibitem{react}
	React documentation [Электронный ресурс]. \mdash Режим доступа: \url{https://react.dev/}. \mdash Дата доступа: 23.09.2026.

	\bibitem{sklearnmlp}
	MLPClassifier [Электронный ресурс]. \mdash Режим доступа: \url{https://scikit-learn.org/stable/modules/generated/sklearn.neural_network.MLPClassifier.html}. \mdash Дата доступа: 23.09.2026.
\end{thebibliography}
